% Generated from locale/tr/content/proof-theory/sequent-calculus/interpretation-rules.tex; do not hand-edit.


\olfileid[tr]{pt}{seq}{irl}

\olsection{Kuralların Yorumlanması}

Bir sekantın yorumunu anımsayalım: $\Gamma \Sequent \Delta$ sekantı,
ya en az bir $!A \in \Gamma$ yanlışsa ya da en az bir $!A \in \Delta$
doğruysa bir !!a{structure}~$\Struct{M}$ içinde doğrudur. \Log{G3c}'nin
mantıksal kurallarını, ana !!{formula}{}ünün doğruluk koşullarını
kodluyor gibi okuyabiliriz. \RightR{\lif}'i ele alalım. Ana
!!{formula}{}ü $!A \lif !B$'dir. $!A$ yanlışsa ya da $!B$ doğruysa bu
formül doğrudur. Dolayısıyla $\ \Sequent !A \lif !B$ sekantı ancak ve
ancak $!A \Sequent !B$ doğruysa doğrudur. Bu sekantların önbileşen ve
artbileşenlerine $\Gamma$, $\Delta$ bağlam !!{formula}s{}ini eklediğimizde
tam olarak \RightR{\lif} kuralının sonucu ile öncülünü elde ederiz.
\LeftR{\lif} için durum benzerdir. $!A \lif !B$, $!A$ doğru ve $!B$
yanlışsa yanlıştır. Öyleyse $!A \lif !B \Sequent \ $ ancak ve ancak hem
$\ \Sequent !A$ hem de $!B \Sequent \ $ doğruysa doğrudur.
\LeftR{\lif}'in sonucu, öncüllerinin ve bağlacıyla kurulan bileşimine denktir. Bu
örüntü, \Log{G3c}'nin kurallarının sağlam olmasını sağlar (bütün öncüller
doğruysa sonuç da doğrudur). Kurallar aynı zamanda tamlığı da sağlar.

Aslında doğruluk fonksiyonlu herhangi bir bağlacın doğruluk koşulları,
bir sekant kümesi olarak bu biçimde betimlenebilir. Bunun nedeni, klasik
mantıktaki her doğruluk fonksiyonunun ve bağlaçlı normal biçimde bir
!!a{formula} ile, yani atomik ve değillenmiş atomik !!{formula}s{}in veya
bağlaçlı bileşenlerinin ve bağlacıyla birleştirilmesiyle betimlenebilmesidir.
Örneğin $!A
\oplus !B$'nin, $!A$ ile $!B$'den tam olarak biri doğru olduğunda, yani
``dışlayıcı veya'' durumunda doğru olduğunu varsayalım. O zaman $!A \oplus
!B$ ancak ve ancak $(!A \lor !B) \land (\lnot !A \lor \lnot !B)$ doğruysa
doğrudur; $(\lnot !A \lor !B) \land (!A \lor \lnot !B)$ doğruysa
yanlıştır. Dolayısıyla $\oplus$ için sekant hesabı kurallarını şöyle
tanıtabilirdik:
\begin{align*}
&
\Axiom$\Gamma \fCenter \Delta, !A, !B$
\Axiom$!A, !B, \Gamma \fCenter \Delta$
\RightLabel{\RightR{\oplus}}
\BinaryInf$\Gamma \fCenter \Delta, !A \oplus !B$
\DisplayProof
\\
& \Axiom$!A, \Gamma \fCenter \Delta, !B$
\Axiom$!B, \Gamma \fCenter \Delta, !A$
\RightLabel{\LeftR{\oplus}}
\BinaryInf$!A \oplus !B, \Gamma \fCenter \Delta$
\DisplayProof
\end{align*}
Bu kurallar da sağlam ve tam olurdu. Sonuçları ancak ve ancak bütün
öncülleri doğruysa doğrudur.

\begin{prob}
$!A \downarrow !B$'nin ancak ve ancak ne $!A$ ne de $!B$ doğruysa doğru
olduğunu varsayın (``NOR''). Bunun \Log{G3c} kuralları nasıl olurdu?
($!A \downarrow !B$ doğru olduğunda ve yalnız o zaman aynı anda doğru olan
iki sekant bulun. Bunlar \RightR{\downarrow} kuralını verir. Ardından $!A
\downarrow !B$ yanlış olduğunda ve yalnız o zaman doğru olan bir sekant
bulun. Bu da \LeftR{\downarrow} kuralını verir.)
\end{prob}

\Log{G3c}'de her bağlaç ve niceleyicinin tam iki kuralı vardır: bir sol
kural ve bir sağ kural. Buna karşılık \Log{G1c}'de iki \LeftR{\land} ve
iki \RightR{\lor} kuralı vardır. Bunun nedeni, \Log{G1c}'nin örnek aldığı
Gentzen'in özgün sistemi~\Log{LK}'nin önemli bir klasik olmayan mantık,
yani sezgici mantık için bir \Log{LJ} çeşitlemesinin bulunmasıdır. Sezgici
mantık için sağlam ve tam bir sistem elde etmek amacıyla \Log{LJ}'deki her
sekantın artbileşeni en fazla bir !!{formula} içerecek biçimde
sınırlandırılmıştır. Bu, öncülünün artbileşeninde hem $!A$ hem de $!B$
bulunan \Log{G3c}'nin \RightR{\lor} kuralını kullanmayı olanaksız kılar.
Bu nedenle Gentzen \Log{LJ} için bir çift \RightR{\lor} kuralı kullandı ve
aynısını \Log{LK} için de kullandı. Benzer biçimde sezgici sürüm~\Log{G1i},
bütün artbileşenleri en fazla bir !!{formula} ile sınırlandırılmış
\Log{G1c}'den ibarettir. (\Log{G3c}'nin de bir sezgici sürümü vardır; ancak
bazı kuralları \Log{G3c}'deki karşılıklarından ayrılır. Örneğin o da bir
çift \RightR{\lor} kuralı kullanır.)

\Log{G1c}'nin yapısal kurallarının sağlam olduğunu görmek kolaydır. Örneğin
$\Gamma \Sequent \Delta$ sağda doğru ya da solda yanlış bir !!{formula}
içeriyorsa $\Gamma \Sequent \Delta, !A$ da aynı !!{formula}{}ü içerir; $!A$'nın
doğru ya da yanlış olması önemsizdir. Dolayısıyla \RightR{\Weakening}'in
öncülü doğruysa sonucu da doğrudur. Burada tersinin geçerli olmadığına
dikkat edin. Sonuç $!A$ doğru olduğu için doğru olabilir; bu durumda
öncülün doğru olması güvence altında değildir.

Peki yapısal kurallara neden gerek vardır? Örneğin $\ \Sequent !A \lor
\lnot !A$'yı !!{prove}{}mak için büzülme (\RightR{\Contraction},
\LeftR{\Contraction}) gerekir. $!A \Sequent !A$ ile başlarsanız önce
\RightR{\lnot} ile $\Sequent !A, \lnot !A$ elde edilir. Sonra
\RightR{\lor}'un iki sürümü art arda kullanılabilir: biri $\lnot !A$'dan,
öteki $!A$'dan $!A \lor \lnot !A$ elde etmek için. Ortaya çıkan sekant
$\ \Sequent !A \lor \lnot !A, !A \lor \lnot !A$'dır. Öyleyse
\Log{G1c}'de $\Sequent !A \lor \lnot !A$ elde etmek için bir !!a{proof}
içindeki aynı !!{formula}{}ün iki geçişini ``büzebilmek'' gerekir:
\[
\Axiom$!A \fCenter !A$
\RightLabel{\RightR{\lnot}}
\UnaryInf$\fCenter !A, \lnot !A$
\RightLabel{\RightR{\lor}}
\UnaryInf$\fCenter !A, !A \lor \lnot !A$
\RightLabel{\RightR{\lor}}
\UnaryInf$\fCenter !A \lor \lnot !A, !A \lor \lnot !A$
\RightLabel{\RightR{\Contraction}}
\UnaryInf$\fCenter !A \lor \lnot !A$
\DisplayProof
\]
Gerçekte zayıflatma ya da büzülme olmadan \Log{G1c}'de !!{prove}{}yamayacağınız
geçerli sekantlar vardır. Örneğin \Log{G1c}'de $!A \Sequent
!B \lif !A$ sekantının bir !!a{proof}{}ı \RightR{\lif} ile bitmelidir
(çünkü yalnız $!B \lif !A$ ana formül olabilir) ve bunun öncül olarak
$!B, !A \Sequent !A$'ya gereksinimi vardır. Bu da $!A \Sequent !A$
aksiyomundan \LeftR{\Weakening} ile elde edilir:
\[
\Axiom$!A \fCenter !A$
\RightLabel{\LeftR{\Weakening}}
\UnaryInf$!B, !A \fCenter !A$
\RightLabel{\RightR{\lif}}
\UnaryInf$!A \fCenter !B \lif !A$
\DisplayProof
\]

\Log{G3c}'de büzülme ve zayıflatma kurallarına hiç gerek yoktur.
Zayıflatma aksiyomların içine yerleştirilmiştir. İki öncüllü kurallarda
$\Gamma$, $\Delta$ bağlamının iki kopyası birleştirilerek büzülmeden büyük
ölçüde kaçınılabilir. Hem \Log{G1c} hem de \Log{G3c} böyle ``bağlam
paylaşan'' kurallara sahiptir. Tek öncüllü kurallarda büzülme yalnız
\RightR{\lor}, \LeftR{\land}, \RightR{\lexists} ve \LeftR{\lforall}
için gerekir. Bu kuralların \Log{G3c}'deki sürümleri büzülmeyi tümüyle
gereksiz kılar. İki öncüllü kuralların bağımsız bağlamlara sahip olduğu bir
\Log{G2c} sistemi vardır (bkz. \olref{tab:G2c}). \Log{G2c}'de büzülme,
\Log{G1c}'dekinden bile daha önemlidir.

\allmodulessubfile{OLP-0706}

